\begin{ledgroupsized}[r]{120mm}
\footnotesize
\pstart
\noindent\textbf{Überlieferung:}
\pend
\end{ledgroupsized}
\begin{ledgroupsized}[r]{114mm}
\footnotesize
\pstart
\parindent -6mm
\makebox[6mm][l]{\textit{L}}%
Konzept:
LH~XXXV~10, 14~Bl.~6. 
Ein Blatt~2\textsuperscript{o};
\href{\Scan{LH_35_10_14_006wz}}%
{Wasserzeichen} in der Blattmitte (Papier aus Osterode im Harz);
linker Rand beschnitten; horizontal gefaltet, im Falz Papierbruch ohne Textverlust.
Eine Seiten auf Bl.~6~r\textsuperscript{o}; Bl.~6~v\textsuperscript{o} leer.
Der Träger bildete ursprünglich eine Einheit mit LH~XXXV~10, 14~Bl.~5 (Träger von N.\;\ref{60114}).
Am oberen Rand von Bl.~6~r\textsuperscript{o}, in Rot, der Vermerk von unklarer Hand \textit{VI}
(vgl. die Träger der benachbarten Signaturen LH~XXXV~10, 14 Bl.~2 bis 13, u.\,a. 
N.\;\ref{60117},
\ref{60118},
\ref{60115},
\ref{60112},
\ref{60116},
\ref{60114},
\ref{60110},
\ref{41824}
und \ref{60111}
in diesem Band).
\pend
\end{ledgroupsized}
\vspace{5mm}
\begin{ledgroup}
\footnotesize
\pstart
\noindent
\textbf{Datierungsgründe:}
Der Träger des vorliegenden Stücks N.\;\ref{60113} hing ursprünglich mit dem Quartblatt LH~XXXV~10, 14~Bl.~5 zusammen, auf dem das Konzept "Perinde est ac si grave impellatur in normalibus ad plana" (N.\;\ref{60114}) überliefert ist. 
Zusammen bilden die Blätter etwa vier Fünftel eines Bogens, dessen verbleibenden Teile bislang nicht identifiziert werden konnten.
Der Träger von N.\;\ref{60113} enthält ein Wasserzeichen mit dem Motiv "Wappen der Stadt Osterode\indexO{Osterode} (Harz)"; N.\;\ref{60114} enthält die Gegenmarke; das Zeichen ist für den Zeitraum 1690 bis 1694 belegt, also für die Zeit unmittelbar nach der italienischen Reise (wobei Leibniz erst in der zweiten Junihälfte 1690 in Hannover eintraf).
\pend
\pstart
Während der Terminus post quem für die Entstehung von N.\;\ref{60113} als sicher gelten darf,
ist der vorgeschlagene Endpunkt lediglich als wahrscheinlich anzusehen.
Die Angabe beruht auf der Feststellung, dass Leibniz sich im Jahr 1690 besonders intensiv mit dem hier erörterten statischen Problem beschäftigte und diese und eng verwandte Fragen durch konsequente Heranziehung seiner Analysis ("Methodus de Maximis et Minimis", hier auf \refnote{FnC_35_10_14_006-maximis}) zu lösen suchte (siehe dazu die auf 1690 datierbaren Stücke N.\;\ref{60112}, N.\;\ref{60220}, und N.\;\ref{41823}).
\pend 
\end{ledgroup}
\vspace*{8mm}
\selectlanguage{latin}
\pstart
\noindent
\href{\Scan{LH_35_10_14_006r}}%
{\lbrack6~r\textsuperscript{o}\rbrack}\ 
\pend
\nobreak
\begin{edfigure}[0em][1em]
\includegraphics[scale=.8]{%
gesamttex/edit_VIII-4/images/LH_35_10_14_006_d1_006r.pdf}
\edcaption{\edtext{\bracks{\textit{Fig.~1, gestrichen}}}{%
\lemma{\bracks{\textit{Fig.~1}}}%
\Cfootnote[nonum]{%
Ein gestrichener Entwurf zum Diagramm wird nicht wiedergegeben.}}}
\end{edfigure}
\begin{edfigure}[1em][1em]
\includegraphics[scale=1]{%
gesamttex/edit_VIII-4/images/LH_35_10_14_006_d2_006r.pdf}
\edcaption[-0.2em]{%
\edtext{\bracks{\textit{Fig.~2}}%
\hspace{10mm}}{%
\lemma{}%
\Afootnote[nonum]{\textit{Neben \bracks{\textit{Fig.~2}}}:
Nihil debet ad r\angles{em} fa\angles{c}ere quae sit figura corporis \textit{ABC} modo normales ad plana incurrant in centrum gravitatis\indexS{centrum gravitatis}}}}
\end{edfigure}
\pstart
\noindent
Sit gravis sphaera \textit{ABC} sustentata\indexS{sphaera a duobus planis sustentata} a duobus planis\indexS{planum inclinatum} \textit{DB}, \textit{EC}, quaeritur quam cuique vim imprimat, seu quae sint pondera\indexS{pondera quibus plana sustinentur} \textit{F}, \textit{G} quibus plana\indexS{planum inclinatum} sustinentur. Pono autem 
\edtext{angulos \textit{BDF} et \textit{CEG} esse rectos et \textit{DE}}{%
\lemma{angulos \textit{BDF}}%
\Bfootnote{%
\textit{(1)}~esse rectos, %
\textit{(2)}~et \textit{CEG} esse rectos %
\textit{(a)}~.~In rect %
\textit{(b)}~et \textit{DE}~\textit{L}}}
esse horizontalem. Ex \textit{A, F, G}, in \textit{DE} agantur 
\edtext{normales, \textit{AH},}{%
\lemma{normales,}%
\Bfootnote{%
\textit{(1)}~\textit{H}, \textit{K}, \textit{L} %
\textit{(2)}~\textit{AH},~\textit{L}}}
\textit{FK}, \textit{GL}. 
Quaeritur situs talis, ut centrum gravitatis\indexS{centrum gravitatis} sit omnium 
\edtext{possibilium infime\indexS{situs centri gravitatis infimus}}{%
\lemma{possibilium}%
\Bfootnote{%
\textit{(1)}~minimum %
\textit{(2)}~infime~\textit{L}}}
situm. Sint pondera \textit{a}, \textit{f}, \textit{g}. 
Erit $AH\cdot a+FK\cdot f+GL\cdot l$ divis. per $a+f+l$ distantia centri gravitatis\indexS{centrum gravitatis} a \textit{DE}, 
quae debet 
\edtext{$\stackrel{(1)}{}$}{\lemma{}\Bfootnote{$\stackrel{(1)}{}$ \textit{erg.~L}}}
esse omnium possibilium maxima.\indexS{distantia centri gravitatis maxima} 
Assumamus \textit{DH} et \textit{AH}, inde caeteri situs 
\edtext{determinabuntur. $DA\!\stackrel{(2)}{=}$}{%
\lemma{determinabuntur.}%
\Bfootnote{%
\textit{(1)}~Datur
\textit{(2)}~$DA\!\stackrel{(2)}{=}$~\textit{L}}}%
$\;\sqrt{DH^2+AH^2}$.
Datur \textit{AB} radius datae sphaerae; cujus quadratum detractum a quadrato \textit{DA}, 
relinquet quadratum \textit{DB}, ergo fiet $DB\stackrel{(3)}{=}\sqrt{DH^2+AH^2-AB^2}$.
Similiter $CE\stackrel{(4)}{=}\sqrt{EH^2+AH^2-AB^2}$. 
Est autem $EH\stackrel{(5)}{=}DE-DH$\bracks{,} 
$DB:BM:MD\stackrel{(6)}{\squaredots}AB:\stackrel{DH-DM}{MH}:MB-AH$. 
Seu $BM:MD\stackrel{(7)}{=}DH-DM,:,MB-AH$ et fiet 
$BM^2-AH\cdot BM\stackrel{(8)}{=}DH\cdot DM\,$%
\edtext{$-\,DM^2$; seu}{%
\lemma{$-\,DM^2$;}%
\Bfootnote{%
\textit{(1)}~seu $DM^2+DM^2\!=$
\textit{(2)}~seu~\textit{L}}}
$AH\cdot BM+DH\cdot DM\stackrel{(9)}{=}BM^2+DM^2$%
\edtext{$\;\stackrel{(10)}{=}DH^2\,$}{%
\lemma{$\;\stackrel{(10)}{=}\,$}%
\Bfootnote{%
\textit{(1)}~$DB^2\stackrel{(11)}{=}$
\textit{(2)}~$DH^2$~\textit{L}}}%
$+\,AH^2-AB^2\stackrel{(11)}{=}AH\cdot BM+DH\cdot DM$ 
quarum duarum aequationum 10 et 11. ope dantur \textit{BM} et \textit{DM}. 
Similiter loco aeqq. 10 et 11. scribemus $CN^2+EN^2\stackrel{(12)}{=}EH^2+AH^2-AB^2\stackrel{(13)}{=}AH\cdot CN+EH\cdot EN$. 
Ita similiter dantur \textit{CN} et \textit{EN} tantumque ad aeqq.~10. 11. 12. 13. 
\edtext{adjungenda ae\bracks{q.}~5. Ponimus}{%
\lemma{adjungenda}%
\Bfootnote{\bars{aeqq. \textit{ändert Hrsg.}} 5. %
\textit{(1)}~Postremo %
\textit{(2)}~Ponimus~\textit{L}}}
autem $DF\stackrel{(14)}{=}EG$, datas. Estque $FK\stackrel{(15)}{=}DF,BM-AH,:AB$ et similiter $GL\stackrel{(16)}{=}DF,CN-AH,:AB$. 
\edlabel{35_10_14_006r_21a}%
\edtext{}{%
\xxref{35_10_14_006r_21a}{35_10_14_006r_21b}%
\lemma{Ex aeq.~1. \bracks{...} $d.AH$}%
\Cfootnote{%
Hier und in der Folge dient das Zeichen \textit{d} bzw. \textit{d.} als Differential-Symbol.}}%
Ex aeq.~1. fit $a\,d.AH + f\,d.FK + l\,d.GL\,$%
\edtext{$\stackrel{(17)}{=}0$ et}{%
\lemma{$\stackrel{(17)}{=}0$}%
\Bfootnote{%
\textit{(1)}~et ex 15 fit %
\textit{(a)}~$DF$ %
\textit{(b)}~$AB\,d\,FK=$
\textit{(aa)}~\angles{$AB$}
\textit{(bb)}~$\overline{BM-AH}\,d\,DF:$
\textit{(aaa)}~$AB$
\textit{(bbb)}~$AB+\overline{DF\, d\, BM}-DF.\overline{d\,AH}.$ %
\textit{(2)}~et~\textit{L}}} 
ex 15 fit $AB.d.FK\stackrel{(18)}{=}DF\,d.BM-DF\,d.AH$.%
\edlabel{35_10_14_006r_21b}
Ex 11 est $DM\stackrel{(19)}{=}DH^2+AH^2-AB^2\,-$%
\edtext{$AH\cdot BM,:DH$ et quadrando 19 fit}{%
\lemma{$AH\cdot BM,:DH$}%
\Bfootnote{%
\textit{(1)}~et $DM^2=$
\textit{(2)}~et quadrando 19 fit~\textit{L}}}
$\ovalbox{DH^4}+AH^4+AB^4+AH^2\cdot BM^2+\ovalbox{2}\cdot DH^2\cdot AH^2+2\cdot AB^2\cdot AH\cdot BM-\ovalbox{2}\cdot AB^2\cdot DH^2-2AB^2\cdot AH^2-2AH\cdot BM\cdot DH^2-2AH^3\cdot BM=\ovalbox{DH^2}\ovalbox{+AH^2\cdot DH^4}\ovalbox{-AB^2\cdot DH^2}$%
\edlabel{35_10_14_006r_1a}%
\edtext{}{%
\xxref{35_10_14_006r_1a}{35_10_14_006r_1b}%
\lemma{$-BM^2\cdot DH^2$. \bracks{/}}%
\Bfootnote{%
\textit{(1)}~Vel si placet brevius procedere $DM=DB^2-AH\cdot BM,:DH. \ovalbox{11}$ (per aeq. 13.) ergo per 9 fit $DB^4+AH^2\cdot BM^2-2DB^2\cdot AH\cdot BM-DH^2\cdot BM^2=DH^2\cdot DB^2=DB^2\cdot DH^2-$ %
\textit{(2)}~$DM^2=DB^2-BM^2=DB^4+AH^2\cdot BM^2-2DB^2\cdot AH.BM,:DH^2$. Ergo $DB^2\cdot DH^2-BM^2\cdot DH^2=DB^4+AH^2\cdot BM^2-2.DB^2\cdot AH\cdot BM$\enskip $DB^2+A^2=DH^2+AH^2$. Ergo $-DB^2+DH^2=-AH^2+AB^2.$ Ergo fit $DB^2\cdot\overline{AB^2-AH^2}=\overline{AH^2+DH^2}\; BM^2-2\cdot DB^2\cdot AH\cdot BM$ seu $DM^2\cdot AB^2-DM^2\cdot AH^2+BM^2\cdot AB^2$\enskip $BM^2\cdot AH^2=AH^2\cdot BM^2+DM^2\cdot BM^2=2\cdot AH\cdot BM^3-2AH\cdot BM\cdot DM^2$\enskip $AH^4+AB^4-2AH^2\cdot AB^2$\bracks{,}
\textit{(a)}~$dB$ %
\textit{(b)}~$BM\; d\; BM+DM\cdot d\; DM=DH\;d \; DH+AH\;d \; AH$ %
\textit{(aa)}~$=AH\;d \; BM+BM\;d \; AH+DH\; d$ %
\textit{(aaa)}~$\ovalbox{+DM.DH}$ %
\textit{(bbb)}~$DM\; d\; DH$. %
\textit{(bb)}~et $2DH\; d\; DH+2AH\; d\; AH=AH\; d\; BM+BM\; d\; AH+DH\; d\; DM+DM\; d\; DH$ %
\textit{(3)}~Sed~\textit{L}}}%
$-\,BM^2\cdot DH^2$.
\pend
\pstart
Sed\edlabel{35_10_14_006r_1b} 
possumus tamen et sic procedere, ut consideremus situm ipsius \textit{EC} qualiscunque 
sit tanquam datum, inde ponamus decurrere \textit{A}, donec \textit{DB} deprimatur, 
ut amplius descendi non possit\bracks{,} et fiet $a\; d.AH+f\;d.FK=0$. %
\pend
\pstart
Et similiter fiet $a\; d.AH+l\;d.GL=0$.
\pend
\pstart
Ergo $f\; d.FK=l\; d.GL$.\enskip\ 
$d.FK=-\overline{a:f}\; d.AH$.\enskip\ 
$d.GL=-\overline{a:l}\; d.AH$.
\pend
\pstart
Sed haec non sint conciliabilia cum aeq.~17. quae tamen et ipsa videbatur esse debere vera.
Nam fieret: $a\; d.AH-a\cdot d.AH+a\; d.HA=0$ quod \edlabel{35_10_14_006r_2a}absurdum.
\edtext{}{\xxref{35_10_14_006r_2a}{35_10_14_006r_2b}\lemma{absurdum.}%
\Bfootnote{%
\textit{(1)}~Jam video in quo peccatum in aeq.~17. nempe quod %
\textit{(2)}~In eo etiam peccatum %
\textit{(3)}~Hinc~\textit{L}}}
\pend
\pstart
Hinc\edlabel{35_10_14_006r_2b} 
perficitur
\edtext{Methodus de Maximis et Minimis.\indexS{Methodus de Maximis et Minimis}}{%
\applabel{FnC_35_10_14_006-maximis}%
\lemma{Methodus de Maximis et Minimis}%
\Cfootnote{Siehe \textsc{G.\,W.~Leibniz}, \cite{02077}\textit{Nova methodus pro maximis et minimis}, \cite{01023}\textit{AE} Oktober 1684, S.\;467--473.}}
\edtext{Nempe cum casus}{%
\lemma{Nempe}%
\Bfootnote{%
\textit{(1)}~ex casibus %
\textit{(2)}~cum casus~\textit{L}}}
sunt infinities infiniti\indexS{casus infinities infiniti} et
\edtext{ex iis maximum}{%
\lemma{ex iis}%
\Bfootnote{%
\textit{(1)}~maximus
\textit{(2)}~maximum~\textit{L}}}
vel minimum quaeritur possunt alternatim omnes assumi constantes
\edtext{praeter duas.}{%
\lemma{praeter}%
\Bfootnote{%
\textit{(1)}~unum. %
\textit{(2)}~duas.~\textit{L}}} 
Nec licet omnes 
\edtext{assumere pro indeterminatis.}{%
\lemma{}%
\Bfootnote{%
\textit{(1)}~indeterminatas %
\textit{(2)}~pro indeterminatis~\textit{L}}} 
\pend
\pstart
Applicandum hoc ad maximam alicujus superficiei\indexS{superficies concava} concavae latitudinem.\indexS{latitudo superficiei concava maxima} 
Ubi etiam non licet omnes simul assumere ut indeterminatas.
\pend
\clearpage
\pstart
$-AB\;\overline{a:f}\;d.AH=DF\;d.BM-DF\cdot d.AH$. Ergo $d.BM:d.AH=f.DF-a\; AB,:,DF$.
Possumus facere $DB=DF$, ita ut \textit{EG} non sit aeq. \textit{DF} sed \textit{EC}, et fiet ni fallor $d.MB$ ut \textit{DM}.
\edtext{Ergo prodibunt}{%
\lemma{Ergo}%
\Bfootnote{%
\textit{(1)}~fiet
\textit{(2)}~prodibunt~\textit{L}}}
ea quae alias.
\pend
\pstart 
Videndum an non eodem calculo res 
\edtext{demonstretur etsi \textit{A} in quo}{%
\lemma{demonstretur}%
\Bfootnote{%
\textit{(1)}~etsi %
\textit{(2)}~etsi \textit{A} %
\textit{(a)}~Anguli %
\textit{(b)}~ipsae \textit{B} %
\textit{(c)}~in quo~\textit{L}}}
concurrunt normales ad plana\indexS{normales ad plana non in centro concurrentes} non sit centrum gravitatis\indexS{centrum gravitatis} globi.
\pend
